% Generated draft: requires mathematical and source review.
\documentclass[11pt]{article}
\usepackage[margin=25mm]{geometry}
\usepackage{fontspec}
\setmainfont{Latin Modern Roman}
\usepackage{amsmath,amssymb,mathtools,unicode-math}
\setmathfont{Latin Modern Math}
\usepackage{enumitem,needspace,xurl,hyperref}
\setlength{\parindent}{0pt}
\setlength{\parskip}{5pt}
\setlength{\emergencystretch}{3em}
\newcommand{\N}{\mathbb{N}}
\newcommand{\Z}{\mathbb{Z}}
\newcommand{\Q}{\mathbb{Q}}
\newcommand{\R}{\mathbb{R}}
\newcommand{\C}{\mathbb{C}}
\newcommand{\F}{\mathbb{F}}
\newcommand{\A}{\mathbb{A}}
\newcommand{\PP}{\mathbb{P}}
\newcommand{\E}{\mathbb{E}}
\newcommand{\eps}{\varepsilon}
\newcommand{\dd}{\,\mathrm{d}}
\begin{document}
\section*{UnsolvedMath ID: 701144}
\textbf{Source problem code:} catalogue entry
\subsection{Definitions and mathematical statement}
Let \(\mathbb M\) denote the Monster finite simple group. A spin manifold is a smooth oriented manifold \(M\) equipped with a spin structure, namely a lift of its oriented orthonormal frame bundle through the double covering \(\operatorname{Spin}(n)\to\operatorname{SO}(n)\), where \(n=\dim_{\R}M\).

The source asks whether there exists a spin manifold \(M\) such that \(\mathbb M\) acts on its loop space. For a specified interpretation \(\mathcal L(M)\) of the loop space, an action consists of a map
\[
\alpha:\mathbb M\times\mathcal L(M)\longrightarrow\mathcal L(M)
\]
satisfying
\[
\alpha(1,\gamma)=\gamma,\qquad
\alpha(gh,\gamma)=\alpha\bigl(g,\alpha(h,\gamma)\bigr)
\quad(g,h\in\mathbb M,\ \gamma\in\mathcal L(M)).
\]
Thus the existence question has the schematic form
\[
\exists\, M\text{ a spin manifold}\quad
\exists\,\alpha\text{ an action of }\mathbb M\text{ on }\mathcal L(M).
\]
The source does not specify whether \(\mathcal L(M)\) is the free loop space, for example \(C^\infty(S^1,M)\), or a based loop space \(\Omega(M,m_0)\), which requires a chosen point \(m_0\in M\). It also does not specify loop regularity, the topology or smooth structure on the loop space, regularity of the action, faithfulness, nontriviality, or structures that the action must preserve. The displayed formulation therefore records the question but does not resolve these missing requirements.

The source further suggests, tentatively, a \(24\)-dimensional manifold and asks whether it could be hyper-K\"{a}hler and noncompact. A hyper-K\"{a}hler structure means a Riemannian metric with parallel orthogonal complex structures \(I,J,K\) satisfying \(I^2=J^2=K^2=IJK=-\mathrm{id}\). These suggestions are not stated as mandatory hypotheses, and the source does not specify which of them should be imposed together or whether the suggested dimension is real or complex. Connectedness, compactness, and boundary conditions for the primary existence question are also unspecified.
\subsection{Short English statement}
J. McKay asks whether the Monster can act on the loop space of a spin manifold, tentatively suggesting dimension \(24\), a hyper-K\"{a}hler structure, and noncompactness. The intended loop space and the required properties of the action are unspecified, so a precise substantive existence claim requires further clarification.
\subsection{Sources}
\begin{itemize}
\item Local problem registry, supplied complete source record: permanent UnsolvedMath ID \(701144\), exact source code \texttt{catalogue entry}, title \emph{Kourovka Notebook Problem 15.56}. The registry labels the problem \texttt{open}; this is a supplied status, not an independently verified mathematical or literature assessment.
\item J. McKay, Problem 15.56, \emph{The Kourovka Notebook (Unsolved Problems in Group Theory)}, 15th issue (2002). The supplied background also identifies the 21st issue edition (2026), \texttt{arXiv:1401.0300v48}. The collection description names E. I. Khukhro and V. D. Mazurov as editors. These bibliographic details are supplied attribution and have not been independently verified here.
\item F. Hirzebruch, T. Berger, R. Jung, \emph{Manifolds and modular forms}, Vieweg, Braunschweig, 1992. This reference is explicitly suggested in the source statement; no page, theorem, or explanation of its precise relevance is supplied.
\end{itemize}
\end{document}
